\documentclass[a4paper,11pt]{article}
\usepackage{mysty}
\usepackage[margin=24mm]{geometry}
\usepackage[hidelinks,unicode]{hyperref}
\usepackage{booktabs}
\usepackage{array}

\newcommand{\RR}{\mathbb{R}}
\newcommand{\ptil}{\tilde p_t}
\newcommand{\nrm}[1]{\left\lVert #1 \right\rVert}

\title{粒子を動かす規則の比較 --- 学習エネルギー + ULA，SVGD，粒子群最適化，合意形成ダイナミクス}
\author{Belief-Score-Based-Model 実装ノート}
\date{2026-10-14}

\begin{document}
\maketitle

\section{設定と記法}

時刻 $t$ の潜在状態 $x_t \in \RR^d$ の belief を $K$ 個の粒子 $\{x_t^i\}_{i=1}^K$ で表す．
観測 $o_t$，直前の行動 $a_{t-1}$，前時刻の粒子 $\{x_{t-1}^j\}$ が与えられたとき，
本プロジェクトの belief は学習した非正規化対数密度
\begin{equation}
  \log \ptil(x) =
  \begin{cases}
    g_\theta(x; o_0) - \tfrac12 \nrm{x}^2 & (t = 0),\\[2pt]
    f_\theta(x; o_t, a_{t-1}) + \log \dfrac1K \sum_{j=1}^K \exp u_\theta(x, x_{t-1}^j, a_{t-1}) & (t \ge 1)
  \end{cases}
  \label{eq:energy}
\end{equation}
を使う．$f_\theta$ は観測項，$u_\theta$ は遷移項で，第 2 式は
「前粒子を祖先とする混合遷移事前分布 $\times$ 観測尤度」という 1 段の事後分布
$p(x_t \mid o_{1:t}) \propto p(o_t \mid x_t) \frac1K \sum_j p(x_t \mid x_{t-1}^j, a_{t-1})$
の学習版である．以下では $E(x) = -\log \ptil(x)$ をエネルギー，$s(x) = \nabla_x \log \ptil(x)$ をスコアと呼ぶ．
粒子を動かす規則（mover）を替えても式~\eqref{eq:energy} は共通で，変わるのは
「$\ptil$ をどう使って $\{x_{t-1}^i\}$ から $\{x_t^i\}$ を作るか」だけである．

費用の比較のため，$C_f$ を $f_\theta$ を $K$ 個の問合せ点で評価する前向き計算の費用，
$C_u$ を $u_\theta$ を $K \times K$ 対で評価する費用とする．
$K = 16$ では $C_u$ が支配的である（遷移対が 256 対）．

\section{現行: 未調整 Langevin 法（Unadjusted Langevin Algorithm，ULA）}

連鎖の初期値を $x^{(0)}_i$（cold なら $\mathcal N(0, I)$，warm なら $x_{t-1}^i$，初期ステップは観測アンカー
$\mu_\theta(o_0) + \sigma_\theta(o_0) \odot \xi_i$）として，$L$ 回
\begin{equation}
  x_i^{(l+1)} = x_i^{(l)} + \alpha\, s\big(x_i^{(l)}\big) + \sqrt{2 \alpha \tau}\; \xi_i^{(l)}, \qquad \xi_i^{(l)} \sim \mathcal N(0, I)
  \label{eq:ula}
\end{equation}
を繰り返し，$x_t^i = x_i^{(L)}$ とする（$\alpha = 0.02$，$\tau = 1$）．方策はスコア $s(x_t^i)$ も受け取る．

性質: (i) $\alpha \to 0$，$L \to \infty$ で粒子は $\ptil$ に従うサンプルになる（不偏ではなく離散化誤差を持つ）．
(ii) 1 ステップあたりの移動量は $\alpha \nrm{s} + \sqrt{2\alpha\tau}$ 程度で，$L \alpha$ の予算を超える跳びには追従できない．
(iii) 粒子間に相互作用はなく（遷移項を通じて前粒子に依存するだけ），同じ谷に全粒子が落ちることを妨げない．
(iv) 乱数 $\xi$ を外生入力として保存し，方策学習（Proximal Policy Optimization，PPO）の再生で同じ連鎖を復元する．
費用: スコアは自動微分で得るので，1 ステップに前向き $+$ 逆向きで約 $3(C_f + C_u)$，連鎖全体で $3L(C_f + C_u)$ に最終スコア 1 回分が加わる．

\section{Stein 変分勾配降下（Stein Variational Gradient Descent，SVGD）と Stein particle filter}

SVGD（Liu \& Wang 2016）は，粒子集合を Kullback--Leibler（KL）ダイバージェンス $\mathrm{KL}(q \,\|\, \ptil)$ の
再生核ヒルベルト空間（Reproducing Kernel Hilbert Space，RKHS）における最急降下方向
\begin{equation}
  \phi^*(x) = \frac1K \sum_{j=1}^K \Big[ k(x_j, x)\, s(x_j) + \nabla_{x_j} k(x_j, x) \Big]
  \label{eq:svgd}
\end{equation}
で動かす: $x_i \leftarrow x_i + \varepsilon\, \phi^*(x_i)$．
核は通常ガウス核（radial basis function，RBF）
\begin{equation}
  k(x, x') = \exp\!\big(-\nrm{x - x'}^2 / h\big)
\end{equation}
で，帯域 $h$（確率密度の記法 $\exp(-\nrm{x-x'}^2/2\sigma^2)$ では $h = 2\sigma^2$）はメディアン則 $h = \mathrm{med}^2 / \log K$
（$\mathrm{med}$ は粒子間距離の中央値，各粒子から見た核の総重みが約 1 になる尺度）で決める．
第 1 項は近傍粒子のスコアを核で平滑化した引力，第 2 項は
\begin{equation}
  \nabla_{x_j} k(x_j, x) = \frac{2}{h}\,(x - x_j)\, k(x_j, x)
\end{equation}
で粒子どうしを押し離す反発項である．Stein particle filter（Fan, Taghvaei, Chen 2021; Maken, Ramos, Ott 2022）は
予測と更新を 1 段の事後分布（式~\eqref{eq:energy} と同じ形）にまとめ，その上で SVGD を回す粒子フィルタで，
全粒子が等重み，リサンプリングなし，である．

ULA との違い: (i) 乱数を使わない決定論的写像で，粒子集合の関数として再生できる（外生ノイズの保存が不要）．
(ii) 反発項が粒子の重なりを防ぐので，$K$ が小さくても縮退しにくい．
(iii) 固定点は $\ptil$ の近似であり，粒子数有限では分散を過小評価する傾向がある（ノイズ項がないため）．
(iv) 多峰の $\ptil$ では核の帯域より離れた谷の間を粒子は移動しない: 谷に粒子を配る役目は初期値（warm start）が担う．
費用: スコア計算は ULA と同じ $3(C_f + C_u)$ に，核とその勾配の $K^2$ 回の評価（$d$ 次元ベクトルの差とスカラー指数関数で，多層パーセプトロン（multilayer perceptron，MLP）の評価ではない）が加わる．

\section{粒子群最適化（Particle Swarm Optimization，PSO）と粒子フィルタとの併用（PSO-PF）}

\subsection{標準 PSO}
Kennedy \& Eberhart（1995）の規則は，各粒子が速度 $v_i$，自己最良 $p_i$ を持ち，群最良 $g$ と合わせて
\begin{align}
  v_i &\leftarrow w\, v_i + c_1 r_1 \odot (p_i - x_i) + c_2 r_2 \odot (g - x_i), \qquad r_1, r_2 \sim \mathcal U(0,1)^d, \label{eq:pso-v}\\
  x_i &\leftarrow x_i + v_i, \label{eq:pso-x}\\
  p_i &\leftarrow \arg\max\{ J(p_i), J(x_i) \}, \qquad g \leftarrow \arg\max_i J(p_i),
\end{align}
で動く．目的関数 $J$ をここでは $J(x) = \log \ptil(x)$ と取る．勾配は不要で，$J$ の前向き評価だけを使う．
慣性 $w$ と係数 $c_1, c_2$ は Clerc \& Kennedy（2002）の収縮係数
$\chi = 2 / \lvert 2 - \varphi - \sqrt{\varphi^2 - 4\varphi} \rvert$，$\varphi = c_1 + c_2 > 4$
で安定化するのが標準である．

PSO は最適化器であり，群は $\arg\max J$ の 1 点に収束する．
belief の広がりを表す目的には，そのままでは使えない（ULA の $\tau$，SVGD の反発項に当たるものがない）．
一方で，自己最良 $p_i$ は「過去に見た高尤度の状態」の記憶であり，粒子が無情報な領域へ流れても引き戻す力になる．

\subsection{PSO-PF}
重み付き粒子フィルタ（Particle Filter，PF）の退化対策として，リサンプリングの前に PSO の数ステップで粒子を観測尤度の高い領域へ寄せる手法群である．
$J(x) = \log p(o_t \mid x)$（または 1 段事後）とし，式~\eqref{eq:pso-v}--\eqref{eq:pso-x} を数回適用したのち
重み $w_i \propto p(o_t \mid x_i)$ を計算してリサンプリングする．
重みと移動を併用する点で，本プロジェクトの「重みを持たず移動だけで belief を表す」設計とは異なるが，
「観測で粒子の位置を動かす」という思想は同じ側にある．

\section{合意形成ダイナミクス（Consensus-Based Optimization，CBO / Consensus-Based Sampling，CBS）}

Pinnau ら（2017），Carrillo ら（2018）の CBO は PSO の確率微分方程式版で，
エネルギー $E$ に対する重み付き平均（Laplace 原理による $\arg\min E$ の近似）
\begin{equation}
  m_\beta = \frac{\sum_i x_i\, e^{-\beta E(x_i)}}{\sum_i e^{-\beta E(x_i)}}
\end{equation}
へ引き寄せつつ，合意点からの距離に比例した雑音を加える:
\begin{equation}
  dx_i = -\lambda\,(x_i - m_\beta)\, dt + \sigma\, \nrm{x_i - m_\beta}\, dW_i
  \qquad(\text{異方版は } \sigma\, \mathrm{diag}(x_i - m_\beta)\, dW_i).
  \label{eq:cbo}
\end{equation}
Carrillo, Hoffmann, Stuart, Vaes（2022）の CBS はこれをサンプラーに変えたもので，重み付き共分散
$C_\beta = \sum_i (x_i - m_\beta)(x_i - m_\beta)^\top e^{-\beta E(x_i)} / \sum_i e^{-\beta E(x_i)}$ を使い
\begin{equation}
  dx_i = -(x_i - m_\beta)\, dt + \sqrt{2 \lambda^{-1} C_\beta}\; dW_i,
  \qquad \lambda = (1+\beta)^{-1} \text{（サンプリング）},\ \lambda = 1 \text{（最適化）}
  \label{eq:cbs}
\end{equation}
とする．平均場極限は非線形 Fokker--Planck 方程式を持ち，ガウス目標分布に対しては定常状態が目標の平均・共分散に一致し，
収束率が目標によらない（affine 不変）．導関数は不要で，$E$ の前向き評価だけで動く．
PSO から CBO への対応（Grassi \& Pareschi 2021）も平均場の枠で示されている．

ULA との違い: (i) 勾配不要なので $3(C_f + C_u)$ が $(C_f + C_u)$ に落ち，同じ浮動小数点演算数（floating-point operations，FLOPs）で約 3 倍のステップを回せる．
(ii) 雑音の大きさが群の広がり $C_\beta$ で決まるので，群が縮むと雑音も縮む（最適化側）か，
目標の共分散に落ち着く（サンプリング側）．
(iii) 単峰のガウス近似が前提にあり，多峰の $\ptil$ では 1 つの合意点 $m_\beta$ に引き寄せられて峰を 1 つ失う．

\section{種分割（niching）}

複数の峰を同時に保つための PSO の拡張で，粒子を距離で種（species）に分け，種ごとに群最良 $g_s$ を持つ．
Li（2004）の species-based PSO では，適合度順に粒子を並べ，半径 $r_s$ 以内に既存の種の代表がなければ新しい種の代表とする:
\begin{align}
  \text{species}(x_i) &= \arg\min_{s} \big\{ \nrm{x_i - g_s} : \nrm{x_i - g_s} \le r_s \big\}, \\
  v_i &\leftarrow w v_i + c_1 r_1 \odot (p_i - x_i) + c_2 r_2 \odot (g_{\text{species}(x_i)} - x_i).
\end{align}
代表どうしが $r_s$ 以内に近づけば種を併合する．クラスタリング（階層的クラスタリング，k-means）で種を切る変種も多い．
鏡像課題の 4 峰は「4 つの種が別々の谷に居続け，標識を読んだ時点で 3 つが消える」問題であり，
種分割はこの構造を規則として直接書ける．ただし種の数と半径は $\ptil$ の谷の構造を知らずに決める必要があり，
谷が 1 つしかない（エネルギーが多峰を学んでいない）場合には何も生まない．

\section{比較}

\begin{center}
\footnotesize
\begin{tabular}{p{28mm}p{21mm}p{23mm}p{21mm}p{21mm}p{22mm}}
\toprule
 & ULA（現行） & SVGD & PSO & CBS & niching PSO \\
\midrule
必要な情報 & スコア $s$ & スコア $s$ & $\log\ptil$ の値 & $\log\ptil$ の値 & $\log\ptil$ の値 \\
乱数 & あり（保存） & なし & あり（$r_1, r_2$） & あり & あり \\
粒子間相互作用 & なし & 核の引力・反発 & 群最良への引力 & 合意点への引力 & 種内の引力 \\
固定点・定常分布 & $\ptil$（近似） & $\ptil$ の近似（分散過小） & $\arg\max \log\ptil$ & ガウス近似 & 各種の極大 \\
多峰 $\ptil$ の扱い & 初期値次第 & 初期値次第，谷間移動なし & 1 峰に収束 & 1 峰に収束 & 種の数まで保持 \\
縮退の防ぎ方 & 雑音 & 反発項 & なし & なし & 種分割 \\
PPO 再生との相性 & ノイズ保存が必要 & 決定論的で容易 & 乱数保存が必要 & 乱数保存が必要 & 乱数保存＋種の履歴 \\
1 ステップの費用 & $3(C_f + C_u)$ & $3(C_f+C_u) + K^2$ 核 & $C_f + C_u$ & $C_f + C_u$ & $C_f + C_u$ \\
\bottomrule
\end{tabular}
\end{center}

いずれの mover でも，エネルギー $\ptil$ 自体の学習信号（PPO の勾配と，次観測・報酬予測の補助損失）は変わらない．
mover が変えるのは「同じ $\ptil$ から粒子集合をどう作るか」であり，
$\ptil$ が谷を 1 つしか持たなければ，どの mover も多峰の belief を作れない．

\section{本プロジェクトへの組み込み案}

\paragraph{実装} \texttt{EnergyBelief} に \texttt{particle\_mover = ula | svgd | cbs | pso} を加える（既定 \texttt{ula} で既存結果と bit 一致）．
SVGD は式~\eqref{eq:svgd} をスコア計算の直後に挟むだけで，外生ノイズの配管を使わない．
CBS は式~\eqref{eq:cbs} の離散化
$x_i \leftarrow m_\beta + e^{-\Delta t}(x_i - m_\beta) + \sqrt{(1 - e^{-2\Delta t}) \lambda^{-1} C_\beta}\; \xi_i$
を使い，$\xi_i$ は ULA のノイズと同じ経路で保存する．スコアの自動微分が消えるので，
同じ FLOPs 予算（core4 の k16f: 1.1 MFLOP/step）で $L$ を約 3 倍にできる．

\paragraph{実験} まず Mountain Hike と LD2D で，k16f の予算を固定して mover だけを \texttt{ula / svgd / cbs} に替えた 3 条件
$\times$ 5 seed を比較する．測るのは return に加えて，粒子間距離と H$_0$ クラスタ数（縮退・多様性），
粒子からの位置復号 $R^2$（belief の情報量），箱外滞在率（罠耐性）．
多峰性は改訂版の鏡像課題で，\texttt{svgd} と種分割つき \texttt{pso} を比較し，
エネルギー地形の谷の数と標識前後の鏡像ビット復号で判定する．
勾配不要 mover では自動微分が消える分だけ壁時間が短くなるはずなので，$L = 12$ を保ったまま学習時間が半減するかも同時に見る．

\paragraph{期待と限界} SVGD は縮退と再生の配管を同時に改善する最も安全な一手である．
CBS は費用面で有利だが単峰近似を内包する．種分割は多峰性に対する唯一の構造的解だが，
エネルギーが多峰を学ぶ圧力（報酬予測の補助損失，$K = 16$ 以上）がなければ機能しない．

\section*{参考文献}
\begin{itemize}
  \item Kennedy, J., Eberhart, R. (1995). Particle swarm optimization. \emph{ICNN}.
  \item Clerc, M., Kennedy, J. (2002). The particle swarm: explosion, stability, and convergence in a multidimensional complex space. \emph{IEEE TEC}.
  \item Li, X. (2004). Adaptively choosing neighbourhood bests using species in a particle swarm optimizer for multimodal function optimization. \emph{GECCO}.
  \item Liu, Q., Wang, D. (2016). Stein variational gradient descent. \emph{NeurIPS}.
  \item Fan, J., Taghvaei, A., Chen, Y. (2021). Stein particle filtering. arXiv:2106.10568.
  \item Maken, F. A., Ramos, F., Ott, L. (2022). Stein particle filter for nonlinear, non-Gaussian state estimation. arXiv:2202.04213.
  \item Pinnau, R., Totzeck, C., Tse, O., Martin, S. (2017). A consensus-based model for global optimization and its mean-field limit. \emph{M3AS}.
  \item Carrillo, J. A., Choi, Y.-P., Totzeck, C., Tse, O. (2018). An analytical framework for consensus-based global optimization method. \emph{M3AS}.
  \item Carrillo, J. A., Hoffmann, F., Stuart, A. M., Vaes, U. (2022). Consensus-based sampling. \emph{Studies in Applied Mathematics}.
  \item Grassi, S., Pareschi, L. (2021). From particle swarm optimization to consensus based optimization: stochastic modeling and mean-field limit. \emph{M3AS}.
  \item Li, T., Sun, S., Sattar, T. P., Corchado, J. M. (2014). Fight sample degeneracy and impoverishment in particle filters: a review of intelligent approaches. \emph{Expert Systems with Applications}.
\end{itemize}

\end{document}
